\documentclass[10pt,a4paper]{article}
\usepackage[margin=25mm]{geometry}
\usepackage{amsmath,amssymb,lmodern}
\setlength{\emergencystretch}{2em}
\begin{document}
\begin{center}
{\Large Closed graphs and domain density of maximal monotone operators}\\[4pt]
{\large Proof of Conjecture 00000008842}\\[6pt]
ziangni-sys \qquad 4 October 2026
\end{center}

\section*{Operators and maximality}
Let $E$ be a real Banach space with continuous dual $E^*$ and
$A:E\rightrightarrows E^*$ a set-valued operator. Its graph and domain are
\[
G=\{(x,x^*):x^*\in A(x)\},\qquad
D=\{x:\exists x^*\in A(x)\}.
\]
Monotonicity means $\langle x^*-y^*,x-y\rangle\geq0$ for all
$(x,x^*),(y,y^*)\in G$. Maximal monotonicity means monotonicity together
with the absence of a proper monotone graph extension.
Every subset of $E\times E^*$ is the graph of a set-valued operator,
by taking its fibers, so maximality is precisely maximality under
inclusion among monotone subsets. No closedness is included in this
definition. All closures below use the product norm topology.

\section*{A general closedness lemma}
Let $X$ be a topological space and $C(p,q)$ a reflexive, symmetric
relation whose sections $\{p:C(p,q)\}$ are closed for every $q$.
Suppose $G\subseteq X$ is maximal among sets all of whose pairs satisfy $C$.
Then
\[
p\in G\quad\Longleftrightarrow\quad
C(p,q)\text{ for every }q\in G.
\]
The forward implication is pairwise compatibility. For the converse,
if $p$ satisfies all those relations, $G\cup\{p\}$ is compatible:
the pairs in $G$ already are, the new diagonal pair uses reflexivity,
and both orders of each mixed pair use symmetry. Maximality forces $p\in G$.
Consequently
\[
G=\bigcap_{q\in G}\{p:C(p,q)\},
\]
an intersection of closed sets, so $G$ is closed.

Apply this lemma to $X=E\times E^*$ and
$C((x,x^*),(y,y^*))\iff\langle x^*-y^*,x-y\rangle\geq0$.
The relation is reflexive, and it is symmetric because negating both
differences leaves the pairing unchanged. Each section is closed by
continuity of the duality pairing: evaluation is a bounded bilinear map,
with $|x^*(x)|\leq\|x^*\|\|x\|$. Hence a maximal monotone graph is closed.
The same argument applies to operators $H\rightrightarrows H$ on a real
Hilbert space using its continuous inner product.

\section*{Density in the closure}
For any topological space and any subset $D$, the canonical inclusion
\[
\iota:D\longrightarrow\overline D,\qquad \iota(x)=x
\]
has dense range. Indeed, a relative open neighborhood of
$x\in\overline D$ is obtained from an ambient open neighborhood of $x$;
the latter intersects $D$ by the definition of closure.
This proves the domain clause, including when the domain is empty.
It asserts relative density in $\overline D$, not density in the entire
ambient space. Combining it with the closedness lemma proves both
clauses of the source conjecture.

\section*{Formal verification}
Lean proves the general closed-compatibility lemma and the actual
\texttt{DenseRange} of the subtype inclusion. It defines graphs and
monotonicity using continuous-dual evaluation and, separately, inner
products; the two \texttt{conjecture\_8842} theorems prove the full
conjunction in these settings. They hold already in real normed spaces
and inner-product spaces: completeness is unnecessary.
\end{document}
